% Generated by report/tools/tables.py from the result files of the repository. Do not edit.
\begin{table}[tbp]
\centering
\footnotesize
\setlength{\tabcolsep}{3.5pt}
\begin{threeparttable}
\caption[Accord entre lectures à l'aveugle]{Accord entre deux lectures à l'aveugle d'agents IA, sur les lectures brutes, avant toute harmonisation. \textbf{En gras}, le seul critère sous 70~\% d'accord, signalé comme fragile par le script : Y1, mesure de performance, avant que ses règles soient précisées (écarts 17 et 19). Sources : \file{study/agree\_v2.json}, \file{study/audit\_stats.json}.}
\label{tab:accord}
\begin{tabularx}{\linewidth}{@{}l>{\raggedright\arraybackslash}Xrrrrrrrr@{}}
\toprule
 & & \multicolumn{2}{c}{\makecell{G\\{}tirage 2026, v1}} & \multicolumn{2}{c}{\makecell{G2\\{}tirage 2026, v2}} & \multicolumn{2}{c}{\makecell{R2\\{}articles ajoutés}} & \multicolumn{2}{c}{\makecell{R3\\{}trajectoire}} \\
\cmidrule(lr){3-4}\cmidrule(lr){5-6}\cmidrule(lr){7-8}\cmidrule(lr){9-10}
Code & Critère & accord & $\kappa$ & accord & $\kappa$ & accord & $\kappa$ & accord & $\kappa$ \\
\midrule
Q1 & L'article a des expériences & 10/10 &  &  &  &  &  &  &  \\
Q2 & Réseaux de neurones employés & 10/10 & 1,00 &  &  & 10/10 & 1,00 & 10/10 &  \\
Q3 & Tous les hyperparamètres de la méthode & 9/10 & 0,80 &  &  & 9/10 & 0,80 & 10/10 & 1,00 \\
Q4 & Tous les hyperparamètres des références & 10/10 & 1,00 &  &  & 10/10 & 1,00 & 10/10 & 1,00 \\
Q5 & Code en lien & 10/10 & 1,00 &  &  & 10/10 & 1,00 & 10/10 & 1,00 \\
Q6 & Choix des hyperparamètres décrit & 10/10 & 1,00 &  &  & 9/10 & 0,62 & 8/10 & 0,62 \\
Q7 & Évaluation sur conditions tenues à part & 9/10 & 0,78 &  &  & 10/10 & 1,00 & 8/10 & 0,41 \\
Q8 & Test statistique appliqué & 10/10 & 1,00 &  &  & 9/10 & 0,62 & 10/10 & 1,00 \\
X2 & Incertitude rapportée & 10/10 & 1,00 &  &  & 10/10 & 1,00 & 10/10 &  \\
X3 & Seuls les meilleurs essais gardés & 10/10 &  &  &  & 10/10 &  & 9/10 & 0,74 \\
X4 & Références réglées avec le même soin & 9/10 & 0,74 &  &  & 10/10 &  & 10/10 & 1,00 \\
X5 & Famille (llm, classic, theory) & 10/10 & 1,00 &  &  &  &  &  &  \\
X7 & Calcul chiffré & 10/10 & 1,00 &  &  & 10/10 & 1,00 & 10/10 & 1,00 \\
\textbf{Y1} & \textbf{Mesure de performance définie} &  &  & \textbf{6/10} & \textbf{0,31} & \textbf{10/10} & \textbf{1,00} & \textbf{10/10} & \textbf{1,00} \\
Y2 & Barres d'erreur définies &  &  & 5/5 & 1,00 & 9/9 & 1,00 & 10/10 & 1,00 \\
Y3 & Plusieurs tâches (au moins 2) &  &  & 10/10 & 1,00 & 10/10 &  & 10/10 & 1,00 \\
Y4 & Une tâche perdue montrée &  &  & 8/8 & 1,00 & 8/10 & 0,52 & 9/9 & 1,00 \\
Y5 & Origine du code des références &  &  & 8/10 & 0,52 & 9/10 & 0,78 & 8/10 & 0,60 \\
Y6 & Exclusions dites &  &  & 9/10 & 0,80 & 8/10 & 0,41 & 8/10 & 0,52 \\
Y7 & Version des bancs d'essai &  &  & 10/10 & 1,00 & 10/10 & 1,00 & 10/10 & 1,00 \\
Y8 & Test de généralisation &  &  & 10/10 & 1,00 & 10/10 & 1,00 & 8/10 & 0,41 \\
Y9 & Sensibilité à un hyperparamètre &  &  & 10/10 & 1,00 & 10/10 & 1,00 & 10/10 & 1,00 \\
Y10 & Versions des bibliothèques &  &  & 10/10 &  & 10/10 & 1,00 & 10/10 &  \\
Y11 & Type d'incertitude nommé &  &  & 5/5 & 1,00 & 9/9 & 1,00 & 10/10 & 1,00 \\
Y12 & Outil statistique postérieur à 2018 &  &  & 10/10 &  & 10/10 & 1,00 & 10/10 & 1,00 \\
\bottomrule
\end{tabularx}
\begin{tablenotes}[flushleft]\scriptsize
\item[] G : dix articles du premier tirage de 2026 relus avec la grille v1 (127 décisions identiques sur 130, \file{study/audit\_stats.json}) ; G2 : les mêmes dix articles relus avec la grille v2 ; R2 : dix articles tirés au hasard parmi les articles ajoutés (\texttt{random.Random(7)}, écart 12), relus avec les deux grilles ; R3 : dix articles de la trajectoire 2019 à 2024 (\texttt{random.Random(2030)}, écart 26). Accord : décisions identiques sur les articles où les deux lecteurs ont répondu oui ou non (Y3 converti en oui à partir de deux tâches) ; $\kappa$ : kappa de Cohen, vide quand un des deux lecteurs ne varie pas (accord trivial). Une case vide hors $\kappa$ : critère absent de cette double lecture. Tous les lecteurs sont des copies du même modèle (écart 28) : l'accord mesure la constance de ce modèle, pas celle de deux lecteurs indépendants.
\end{tablenotes}
\end{threeparttable}
\end{table}
